\section*{Chapitre \ref{ch:s1:value-quality-and-low-risk} --- Valeur, qualité et faible risque}

\begin{solution}{exo:s1:value-quality-and-low-risk:1}
Le gain au sommet était de 8.77~\$. Une chute de 57.8~\% laisse $9.77 \times 0.422 = \$4.12$, soit une perte de 5.65~\$ sur les 8.77~\$ gagnés~: 64~\%.
\end{solution}

\begin{solution}{exo:s1:value-quality-and-low-risk:2}
Position longue de $1/0.6 = \$1.67$ d'actions à bêta faible et position courte de $1/1.5 = \$0.67$ d'actions à bêta élevé~: chaque côté a un bêta égal à un, et le portefeuille est net long
$1.67 - 0.67 = \$1.00$, financé par un emprunt.
\end{solution}

\begin{solution}{exo:s1:value-quality-and-low-risk:3}
Chaque côté rapporte $0.5 \times 1.75 \times 3~\%$, donc le portefeuille rapporte $2 \times 0.5 \times 1.75 \times 0.03 = 5.25~\%$ par an avant coûts, si la prime
était la seule chose que portent ses actions.
\end{solution}

\begin{solution}{exo:s1:value-quality-and-low-risk:4}
La valeur comptable évolue lentement, donc l'essentiel de la variation mensuelle du \omterm{def:s1:value-quality-and-low-risk:value}{ratio valeur comptable sur cours} (book-to-price) est l'opposé du rendement de l'action~: les titres qui ont baissé deviennent bon marché. Un portefeuille value à cours courant
est donc long des perdants récents et court des gagnants récents, à l'opposé du momentum. Deux stratégies aux primes positives et aux rendements
corrélés négativement se combinent en un portefeuille dont le ratio de Sharpe dépasse celui de chacune d'elles.
\end{solution}

\begin{solution}{exo:s1:value-quality-and-low-risk:5}
Son signal, les bénéfices rapportés à la valeur comptable, baisse quand la valeur comptable augmente, donc il classe les titres à forte valeur comptable (bon marché) en bas~: le portefeuille est court en value, et le marché synthétique rémunère une
prime value. Sa corrélation avec le portefeuille value vaut $-0.60$, et le coefficient d'information de son signal, par rapport à la prime value insérée, vaut $-0.29$.
\end{solution}

\begin{solution}{exo:s1:value-quality-and-low-risk:6}
$0.594 \times 7.2~\% = 4.3~\%$ et $1.604 \times 7.2~\% = 11.5~\%$. Les déciles ont rapporté 6.8~\% et 8.6~\%, soit 2.5 points de plus au bêta faible et 2.9 points de moins au
bêta élevé que ne le prévoit le modèle. La droite de marché des titres est plate.
\end{solution}

\begin{solution}{exo:s1:value-quality-and-low-risk:7}
Le ratio de Sharpe du composite vaut 0.44. Le portefeuille value à cours courant et le portefeuille momentum ont une corrélation de $-0.56$, si bien que leurs signaux
se compensent en partie leur bruit mutuel~: le composite conserve les deux primes avec moins de risque que chacun des deux portefeuilles pris isolément.
\end{solution}

\begin{solution}{exo:s1:value-quality-and-low-risk:8}
Quatorze années de pertes sont compatibles avec une prime qui existe encore~: l'écart-type de HML vaut 10~\% par an, si bien qu'une moyenne de 3.6~\% a une erreur type d'environ
2.7~\% sur quatorze ans, et la value a rapporté 9.4~\% par an depuis 2021. La traversée du désert est un élément à peser, non une preuve~; les meilleures questions sont de savoir si sa
mesure capte encore la value (actifs incorporels), et comment la value s'insère dans le reste du portefeuille.
\end{solution}

\begin{solution}{pb:s1:value-quality-and-low-risk:1}
\begin{enumerate}
\item Fonds propres comptables sur cours~; achat des ratios élevés et vente des ratios faibles. HML vaut la moitié de small value, plus la moitié de big value, moins la moitié de small growth et la moitié de big growth,
reconstitué chaque juin à partir de la valeur comptable du dernier exercice et de la capitalisation boursière de décembre.
\item 3.6~\% par an et 0.35 sur 1963--2026~; 5.7~\% jusqu'en 2006, $-5.3~\%$ sur 2007--2020, 9.4~\% depuis 2021.
\item Partie de 9.77~\$ à la fin de 2006, en baisse de 57.8~\% jusqu'en septembre 2020~; encore 29.6~\% sous le sommet de juillet 2026.
\item Le \omterm{def:s1:value-quality-and-low-risk:value}{ratio valeur comptable sur cours} (book-to-price) calculé avec un ancien cours est un signal value lent~; avec le cours d'aujourd'hui, il s'expose aussi aux rendements récents (vente de momentum).
\item Long sur les sociétés rentables, court sur les sociétés non rentables~; le \omterm{def:s1:value-quality-and-low-risk:quality}{facteur qualité} ajoute la croissance et la sécurité. La rentabilité brute prédit les rendements à peu près aussi bien que le ratio valeur comptable sur marché (book-to-market)~; QMJ
rapporte des rendements significatifs ajustés du risque dans 25 pays.
\item 0.09, $-0.19$ et 0.68 sur 1963--2026~; $-0.14$, $-0.48$ et 0.49 sur 2007--2020.
\item 2.8~\% par an, avec un ratio de Sharpe de 0.51.
\item C'est en réalité une position courte en value, et le marché synthétique rémunère la value mais pas la rentabilité.
\item Les actions à faible risque rapportent davantage, ajusté du risque, que ne le prédit le MEDAF~; BAB est long du bêta faible avec levier et court du bêta élevé sans levier.
\item Le décile le plus bas (bêta 0.59) a rapporté 6.8~\% au-dessus des bons du Trésor, et le plus haut (1.60) 8.6~\%, contre 4.3~\% et 11.5~\% prévus.
\item 4.3~\% par an depuis 1968, ratio de Sharpe de 0.26, bêta $-0.06$.
\item La droite de marché des titres du marché synthétique est pentue (2.4~\% pour un bêta de 0.49, 13.4~\% pour 1.39)~: rien n'y rémunère le bêta faible.
\item \textbf{Résultat nommé.} HML a reculé de 57.8~\% depuis son sommet de décembre 2006 jusqu'en septembre 2020~; sa corrélation avec RMW valait 0.09 sur 1963--2026 et $-0.14$ sur
2007--2020.
\item Timing Fama--French 0.24, annuel au cours d'aujourd'hui 0.21, avec biais d'anticipation $-0.06$, mensuel au cours d'aujourd'hui $-0.07$~: plus le signal intègre
les rendements récents, plus il est court en momentum inséré.
\item Il atteint 0.44, car la value à cours courant et le momentum sont corrélés de $-0.56$.
\item Des cours à jour prévoient mieux le vrai \omterm{def:s1:value-quality-and-low-risk:value}{ratio valeur comptable sur cours} (book-to-price) et ont rapporté des alphas de 3.05--3.78 points par an face à la méthode standard.
\item Faire varier les pondérations des facteurs selon les prévisions de leurs rendements. Pour~: les écarts de valorisation et le prix de la qualité ont prédit les rendements. Contre~: signaux lents,
échantillons courts, et le bon marché de la value sur 2007--2020.
\item Dates de dépôt des fondamentaux, timing des cours, années de traversée du désert dans l'échantillon, et corrélations avec les autres facteurs maintenues fixes.
\item Pondérations quasi fixes entre value, momentum et rentabilité, chacune ajustée par la volatilité, le timing n'étant au mieux qu'une légère inclinaison.
\item Chacune est une prime lente, avec de longues périodes de pertes que ses partenaires ont tendance à compenser.
\end{enumerate}
\end{solution}

\begin{solution}{iq:s1:value-quality-and-low-risk:1}
Chaque juin, les actions sont réparties selon la taille médiane du NYSE et selon les 30\up{e} et 70\up{e}~centiles du NYSE du ratio valeur comptable sur marché (book-to-market), soit la valeur comptable du dernier exercice divisée par la
capitalisation boursière de décembre)~; HML est la moyenne des portefeuilles small et big value moins la moyenne des portefeuilles small et big growth, pondérés en valeur et détenus pendant un
an.
\end{solution}

\begin{solution}{iq:s1:value-quality-and-low-risk:2}
Les investisseurs qui ne peuvent ou ne veulent pas recourir à l'effet de levier s'exposent au marché en achetant des actions à bêta élevé, qu'ils font monter~; les gérants évalués par rapport à un indice les préfèrent aussi. Le bêta élevé
s'accompagne alors d'un alpha faible, et la droite de marché des titres est plus plate que ne le dit le MEDAF.
\end{solution}

\begin{solution}{iq:s1:value-quality-and-low-risk:3}
Que la value ait connu un drawdown long et profond, ce qu'un facteur de volatilité de 10~\% peut connaître par hasard avec une prime positive~; que la croissance, les taux bas et les actifs incorporels
ont dominé la période. Non que la value soit morte~: quatorze ans ne suffisent pas à trancher, et elle a récupéré une partie de la perte depuis 2021.
\end{solution}

\begin{solution}{iq:s1:value-quality-and-low-risk:4}
Leurs rendements sont décorrélés ou corrélés négativement à ceux de la value, si bien que la combinaison conserve les primes avec moins de risque~; la rentabilité distingue aussi
les sociétés bon marché de bonne qualité des sociétés bon marché de mauvaise qualité, ce qui améliore la value.
\end{solution}

\begin{solution}{iq:s1:value-quality-and-low-risk:5}
Valeurs comptables connues seulement à partir de leurs dates de dépôt, et non de leurs clôtures d'exercice~; retraitements~; le cours utilisé (courant ou à la clôture de l'exercice)~; nombre d'actions et fractionnements
entre la date de la valeur comptable et aujourd'hui~; sociétés radiées conservées dans l'historique.
\end{solution}

\begin{solution}{iq:s1:value-quality-and-low-risk:6}
Avec des pondérations égales et des volatilités égales $\sigma$, la moyenne vaut $0.3\sigma$ et la variance $\frac14\sigma^2(2 + 2\rho) = \frac14\sigma^2$, si bien que le ratio de Sharpe
vaut $0.3/0.5 = 0.6$. La meilleure combinaison a $SR^2 = s^\top C^{-1} s = 2 \times 0.09/(1 + \rho) = 0.36$, soit encore 0.6~: par symétrie, les pondérations égales sont
optimales.
\end{solution}
